\input{\string~/Documents/School/"UC Irvine"/ta_template2.0.tex}
\rh{2B}{5.2}
\chead{\today}
\graphicspath{ {\string~/Desktop} }
\usetikzlibrary{arrows}
%============ANSWER KEY TOGGLE===========
\newcommand{\ans}[1]{ \noindent {\color{red} \textbf{Answer:} #1}}
%\renewcommand{\ans}[1]{\vspace{3in}}
%==========================================
\begin{document}
\begin{center}
\textbf{Definite Integrals}
\end{center}

This section comes with very minor definitions, which introduces some new notation:

\begin{defi}[Integrable]
A function $f$ is \emph{integrable} on $[a,b]$ if it is continuous or has at most finitely many jump discontinuities on the interval. This means the definite integral exists.
\end{defi}


\begin{defi}
The \textbf{definite integral} of a function $f(x)$ on the interval $[a,b]$ is:
\[
\int_a^b f(x) dx = \lim_{n \to \infty} \sum_{i = 1}^n f(x_i) \Delta x
\]
where $x_i$ can be taken as in the \emph{left-} or \emph{right-}hand sums. These have the properties:
\begin{description}
\item[Negatives] $\int_a^b f(x) dx = -\int_b^a f(x) dx$
\item[Constants] For real number $c$, $\int_a^b c dx = c(b-a)$ and $\int_a^b c f(x) dx = c \int_a^b  f(x) dx$
\item[Sums \& Differences] For two integrable functions $f,g$ $\int_a^b f\pm g dx = \int_a^b f(x) dx \pm \int_a^b g(x) dx$. Also, 
\[
\int_a^b f(x) dx = \int_a^c f(x) dx + \int_c^b f(x) dx
\]
\item[Comparisons] If $m \leq f(x) \leq M$ for all $a \leq x \leq b$, then $$m(b-a) \leq \int_a^b f(x)dx \leq M(b-a)$$. Also if $f(x) \leq g(x)$ for all $a \leq x \leq b$, then $$\int_a^b f(x)dx \leq \int_a^b g(x) dx$$. 

\end{description}
\end{defi}

\begin{exx}
Consider the following integral:
\[
\int_0^1 7x^2 + 4
\]
\begin{enumerate}[label=\alph*)]
\item Approximate the integral with $n=3$ subintervals.
\item Express the integral as the limit of the Riemann sum
\item Compute the definite integral.
\end{enumerate}
\end{exx}
\ans{The final answer for $(c)$ should be $4 + \frac73$}

\problembox{Compute the following integral as the limit of a Riemann sum:
\[
\int_{-2}^0 x^2 + x dx
\]
}

\ans{Using $\Delta x = \frac{0 + 2}{n}$, $f(x) = x^2 + x$, and the summation formulas we have the following steps:
\begin{alignat*}{2}
\int_{-2}^0 x^2 + x dx &= \lim_{n \to \infty} \sum_{i = 1}^n \lrp{\lrp{-2 + i \cdot \frac2n}^2 + \lrp{-2 + i \cdot \frac2n}} \cdot \frac2n \\
&= \lim_{n \to \infty} \sum_{i = 1}^n \lrb{2 - \frac{6i}{n} + \frac{4i^2}{n^2}} \cdot \frac2n \\
&= \lim_{n \to \infty} \lrb{\lrp{\frac4n \sum_{i = 1}^n 1}  - \lrp{ \frac{12}{n^2} \sum_{i = 1}^n i } +\lrp{ \frac{8}{n^3}\sum_{i = 1}^n i^2}} \\
&= \lim_{n \to \infty} \frac4n \cdot n - \frac{12}{n^2} \cdot \frac{n(n+1)}{2} + \frac{8}{n^3} \cdot \frac{n(n+1)(2n+1)}{6} \\
&= \lim_{n \to \infty} 4 - 6 \cdot \frac{n^2 + n}{n^2} + \frac{4}{3} \frac{2n^3 + 3n^2 + n}{n^3}\\
&= 4 - 6 + \frac{8}{3} \\
&= \boxed{\frac{2}{3}}
\end{alignat*}
}

\problembox{Interpreting the graph of these functions, compute the integrals:
\[
\int_{-3}^3 \sqrt{9-x^2} \mte{ and}  \int_0^{37}x dx
\]
}

\ans{First is a circle of radius 3 and second is a triangle.}




\end{document}